\section{The Word: Two Lectionary Provisions}

The most consequential structural difference between the two books is not a prayer at all. It is that one of them contains the readings and the other does not.

\subsection{A plenary Missal and a Missal with a Lectionary}

The 1962 \emph{Missale Romanum} is a plenary book. Every Mass formulary prints its Epistle and its Gospel in place, between the Collect and the Secret, so that the priest at the altar reads the whole Mass from one volume. The postconciliar Missal prints the prayers and the antiphons and refers the readings to a separate \emph{Ordo lectionum Missæ}, first published in 1969 and issued in a second typical edition in 1981, which the current \emph{Institutio generalis} cites by name.

This is not a matter of binding convenience. The 2021 \emph{Responsa ad dubia} treats it explicitly as one of the reform's substantive achievements: the publication of the Lectionary, it says, ``in addition to overcoming the `plenary' form of the \emph{Missale Romanum} of 1962,'' returns ``to the ancient tradition of individual books corresponding to individual ministries.''\endnote{Congregation for Divine Worship and the Discipline of the Sacraments, \emph{Responsa ad dubia} (4 December 2021), explanatory note to the response on the use of the full biblical text; official English text, registered artifact, checked 25 July 2026.} A plenary Missal presupposes that one man says everything; a book set presupposes that a lector, a psalmist, a deacon and a priest each hold their own book. The lectionary question and the ministries question are the same question seen from two sides, and section~8 returns to it.

\subsection{The two provisions, side by side}

The newer provision's architecture is stated in its own \emph{Prænotanda}: the Sunday and festal order runs over three years and the weekday order over two, each proceeding independently of the other; every Sunday and festal Mass has three readings, from the Old Testament, from an apostle, and from a Gospel; the same texts return only every fourth year; and the governing principles are named as ``harmonious composition'' and ``semicontinuous reading,'' applied differently according to the season.\endnote{\emph{Ordo lectionum Missæ}, editio typica altera (1981), \emph{Prænotanda} nn.~65--69, printed pp.~XXXII--XXXIV: ``Ordo lectionum dominicalis--festivus per triennium explicatur, ferialis autem per biennium'' (n.~65); ``Quælibet Missa tres exhibet lectiones \dots{} cyclus trium annorum proponitur, ita ut iidem textus solummodo quarto quoque anno legantur'' (n.~66); and, for weekdays, an annual Gospel cycle in Ordinary Time with a biennial first reading, Year I in odd and Year II in even years (n.~69). Read in the registered scan of the Latin typical edition; nn.~65, 66 and 69 are separately registered passages in the source library.}

\begin{comparetable}{1962 Missal}{\emph{Ordo lectionum Missæ}}{}
Two readings at Mass: Epistle (or a lesson from the Old Testament, Acts, or the Apocalypse) and Gospel & Three readings on Sundays and solemnities of the Lord: Old Testament, apostle, Gospel; two on weekdays & The Old Testament reading is made a normal Sunday element rather than an occasional one\\
One annual cycle: the same Epistle and Gospel return on the same Sunday each year & Three-year Sunday cycle (A, B, C), so that the same texts return every fourth year; weekday Gospels annual, weekday first readings in a two-year cycle (I in odd years, II in even) & Sunday repetition interval extended from one year to three\\
Gradual, Alleluia with verse, or Tract; Sequence on five days & Responsorial psalm; Alleluia or, in Lent, another acclamation; Sequence, obligatory on Easter and Pentecost & The interlectionary chant is retained and redesigned around congregational response\\
Selection principle not stated in the book & Selection principles stated at length in the Lectionary's own \emph{Prænotanda} nn.~64--77: harmony and semicontinuous reading, the reservation of certain books to certain seasons, criteria of length, difficulty, and omission & The newer provision publishes its own rationale and thereby exposes it to criticism\\
\end{comparetable}

\subsection{Two claims that do not survive}

\textbf{``The old Mass had no Old Testament.''} This is false, and demonstrably so from the 1962 book. Old Testament lessons stand in it throughout: the Lenten and Ember ferias draw on Exodus, Leviticus, Numbers, Deuteronomy, the books of Kings, Esther, Daniel, Jeremias, Ezechiel and the Machabees; Isaias supplies a substantial series; the sapiential books supply the standing lessons of many sanctoral Masses and Commons; the Easter Vigil prophecies are a continuous course of Old Testament reading. A survey of the printed incipit formulas in the artifact's text layer counts on the order of one hundred and eighty \emph{Lectio} headings drawn from Old Testament books, against roughly two hundred from the Pauline corpus.\endnote{Counted programmatically over the embedded text layer of the registered facsimile, 25 July 2026, by matching the printed incipit formulas (\emph{Lectio libri \dots{}}, \emph{Lectio Isaiæ Prophetæ}, \emph{Lectio Epistolæ beati Pauli Apostoli}, and the like). Optical-character-recognition error and repeated formularies mean these totals are approximate; the tallies and their method are recorded as a bounded finding in \texttt{research/scope.md}. The count is offered to establish presence and rough proportion, not as an edition statistic, and the qualitative distribution described in the text was checked against the printed formularies themselves.}

What is true, and worth stating exactly, is where those lessons fall. On the Sundays after Pentecost---the long green season in which most Catholics heard most of their Mass readings---the pre-Gospel reading in the 1962 book is drawn from an apostle without exception. The Old Testament in the older Missal is a matter of Lent, the Ember Days, the Vigil, and the sanctoral; it is not a Sunday element. The reform's change was therefore not from absence to presence but from occasional and seasonal to weekly and structural, and that is a large change stated accurately.

\textbf{``Omitting verses is a postconciliar invention.''} The Lectionary's own \emph{Prænotanda} state the practice and deny the novelty in the same paragraph: ``The tradition of many liturgies, the Roman liturgy by no means excluded, has sometimes been accustomed to omit certain verses in the readings from Scripture. It must further be admitted that such omissions cannot be made lightly, lest the sense of the text or the mind and, as it were, the style of Scripture be mutilated. Nevertheless, with care taken that the sense truly remain whole in its essence, it seemed right, on pastoral grounds, to preserve that tradition in the present Order too.''\endnote{\emph{Ordo lectionum Missæ}, editio typica altera (1981), \emph{Prænotanda} n.~77, printed p.~XXXV: ``Traditio multarum liturgiarum, ipsa liturgia Romana minime exclusa, aliquando consuevit in lectionibus Scripturæ quosdam versiculos omittere. Ultro fatendum est huiusmodi omissiones non leviter fieri posse, ne sensus textus aut mens et quasi stilus Scripturæ mutilentur. Attamen, cauto ut sensus revera in sua essentia integer maneat, visum est, ex ratione pastorali, etiam in præsenti Ordinis prædictam traditionem servare.'' Read in a rendered page image of the registered scan of the Latin typical edition, 26 July 2026; the English here is a working gloss.} The 1962 book bears this out on its own pages: its pericopes are editorially bounded, its Epistles are furnished with liturgical incipits (\emph{Fratres}, \emph{Carissimi}) that replace the original openings, and its selections stop and start where the formulary requires rather than where the biblical author does.

\subsection{The criticism that does survive, in the reform's own words}

Against those two corrections stands a criticism that the Lectionary's \emph{Prænotanda} concede in terms. Immediately before the paragraph on omitted verses, under the heading \emph{De textibus difficilioribus}, stands this: ``For pastoral reasons, in the readings on Sundays and solemnities, biblical texts which are genuinely more difficult are avoided.''\endnote{\emph{Ordo lectionum Missæ}, editio typica altera (1981), \emph{Prænotanda} n.~76, printed p.~XXXV: ``Ex ratione pastorali vitantur in lectionibus, diebus dominicis et sollemnitatibus, textus b[i]blici qui revera difficiliores sunt, sive obiective eo quod altiora problemata litteraria, critica aut exegetica movent, sive etiam, quadamtenus saltem, eo quod a fidelibus difficilius intellegi possunt.'' The registered scan prints \emph{bilici} where the sense requires \emph{biblici}; the defect is recorded and not silently corrected. The article continues: ``Attamen divitias spiritales quorundam textuum non licuit fidelibus abscondi, eo quod ab eis difficilius capiantur, quando hæc difficultas oritur sive ex insufficienti institutione christiana, qua nullus fidelis carere debet, sive ex insufficienti institutione biblica, qua omnis pastor animarum pollere debet''---the spiritual riches of certain texts could not be hidden from the faithful on the ground that they find them harder, when the difficulty arises from inadequate Christian or inadequate biblical formation. Read in a rendered page image, 26 July 2026; the English here is a working gloss.} The sentence is not a polemicist's summary; it is the reform's own statement of a selection criterion, printed in the typical edition's own introduction. Anyone who has been told that the charge of pastorally motivated omission is an invention of the reform's critics is holding a claim that the reform's own introduction refutes.

Two qualifications belong immediately beside that concession, and they belong there because they are also in the text. The first is the paragraph's own limiting clause, quoted in the note: difficulty arising from inadequate catechesis or inadequate clerical formation is expressly disallowed as a reason for omission. The second is the \emph{Institutio generalis}, which, in governing the choice among optional texts, warns that ``care must be taken lest, in the selection of texts of Sacred Scripture, parts of it be permanently excluded.''\endnote{\emph{Institutio Generalis Missalis Romani} (2002), n.~361, artifact p.~55: ``Caveatur tamen ne, in seligendis textibus Scripturæ Sacræ, partes eius permanenter excludantur.''} The reform legislated both a policy of omission and a norm against its abuse. A fair account of the lectionary question quotes both.

\subsection{What the comparison cannot settle}

Two questions arise here that the books cannot answer and this study will not pretend to.

The first is whether more Scripture, heard once every three years, forms a congregation better than less Scripture heard every year. That is not a textual question. It is a question about memory, repetition, and catechesis, and the honest answer is that nobody has measured it. The 1962 provision is short and cyclical: a Catholic who attended Sunday Mass for a decade heard one small set of Epistles and Gospels ten times over, in the same order, on the same days. The postconciliar provision is long and progressive: the same Catholic hears a substantially larger body of Scripture, and hears any given passage about three times. These are different pedagogies with different characteristic virtues, and the choice between them is a prudential judgement that the Church made and may revisit.

The second is what \emph{Traditionis custodes} art.~3 \S 3 does to the older provision in practice. Since the readings must now be proclaimed in the vernacular from a conference-approved translation, and since no vernacular lectionary reproducing the older cycle may be published, the 1962 book's pericope list has been detached from the 1962 book's text: the cycle survives as a schedule of references, read from a modern Bible.\endnote{\emph{Traditionis custodes} art.~3 \S 3 and the \emph{Responsa ad dubia} explanatory note cited above. The practical consequence---an older cycle read in a newer translation---is a legal fact of the current discipline, not a judgement about it.} That is a real and recent change in what ``celebrating with the 1962 Missal'' means, and it is one that neither party in the older argument anticipated.
